\documentclass[12pt]{article}
\usepackage[utf8]{inputenc}
\usepackage[T1]{fontenc}
\usepackage{hyperref}
\title{Orya Language Resources for K-12 Curriculum Development}
\author{Compiled by Claude Sonnet 5 (claude-sonnet-5) for Andr\'as Kornai}
\date{}
\begin{document}
\maketitle

\section{Introduction}
Orya (ISO 639-3 ury, also spelled Oria or Uria) is a Papuan language of Papua province, Indonesia. It is the closest relative of the Tor languages, and together they form the Orya--Tor family. Speakers call themselves Yap Zi and their language Yap Zi Ol. Usher's profile gives about 1,600 speakers in 1991. Dialects include West Orya, East Orya and Yapsi-Taja. The Sabda profile rates the language as threatened.

\section{Primary Educational Resources}
I found no primers, textbooks or school programmes for Orya.

\section{Linguistic and Reference Materials}
\begin{itemize}
\item Overview of the language: \url{https://en.wikipedia.org/wiki/Orya_language}
\item Orya--Tor family overview: \url{https://en.wikipedia.org/wiki/Orya%E2%80%93Tor_languages}
\item New Guinea World page on Orya: \url{https://newguineaworld.linguistik.uzh.ch/families/northwest-new-guinea/foja-range/orya-tor-river/orya}
\item Wiktionary category of Orya words: \url{https://en.wiktionary.org/wiki/Category:Orya_language}
\end{itemize}

\section{Religious and Community Resources}
\begin{itemize}
\item Sabda language profile (code ury): \url{https://bakat.sabda.org/502-ury-001.html}
\end{itemize}

\section{Digital and Audio Resources}
None found.

\section{Curriculum Development Context}
The language is still spoken by adults, but I found no teaching materials. A curriculum would need community-led documentation. Indonesian would be the school language alongside it.
\end{document}